\section{What Is Solid, What Is Partial, What Is Open}
\label{sec:status}

A conceptual map is only useful if it is honest about its own standing. The picture above is
not uniform in confidence, and it helps to sort it plainly. This is a summary at the level of
the whole, not a list of the individual experiments, which belong to a separate, more technical
account.

\paragraph{Solid.}
The mathematics at the base is established and standard: the integers, their symmetries, the
Coxeter construction, the hyperbolic crystals, the golden ratio, the fact that finite-celled
hyperbolic crystals exist only in two, three, and four dimensions. On the substrate, the
following are demonstrated in the testbed: a sharp emergent dimension and geometry, an emergent
law of time and energy with the right three properties at once, a Lorentz-safe substrate from
curvature rather than disorder, the field content with mass from a Higgs, gravity's Newtonian
limit and a conserving Einstein equation, the area law, the arrow and expansion, the continuum
limit and its stability under coarse-graining, the unifying integer base, the structural model
of freedom and choice, computational universality, the recursion of higher vibes with its
emergent macro-rule, the nested selves and the tower, the no-complete-self-storage principle,
integration picking out selves, the subtle-layer urge, waking versus dreaming, pattern
persistence, the dimension selection, the Born-rule derivation, Hawking temperature, the
nonlinear Einstein equation in cosmological form, and the discrete graviton.

\paragraph{Partial.}
Several pieces are real but unfinished. The quantum sector has its pillars in place, with the
Born rule now derived, but a full account is still being assembled. Dark energy and dark matter
have mechanisms with the right qualitative behavior, with the precise numbers still ahead. The
discrete graviton is shown as an operator with the right properties, while the second variation
of the full action on a random sprinkling, where fluctuations are large, remains. The large-N
free-energy crossing is unblocked by a new move but not yet driven to full convergence at the
largest sizes. Life as far-from-equilibrium self-maintenance, and evolution as pattern
selection, are sound in outline rather than fully built.

\paragraph{Open.}
Some things are named and not yet built: the specific content of the Standard Model, including
why there are three generations and the pattern of masses and the particular gauge group,
Hawking's late-time information return in microscopic detail, structure formation and the
primordial spectrum, and the strong-field interior. And one thing is open in a deeper sense than
the rest. The model can build and measure the \emph{structure} of experience, including the
structure of a unified self. Whether that structure is the same as the felt fact of
experiencing, the hard problem at the center of the whole picture, is not something a simulation
can settle from the outside. We treat the experiential base as a starting commitment, and the
shared physical pattern as the part that can be checked, and we keep the difference between the
two clearly in view.

\paragraph{The shape of the claim.}
What the program shows is that a strict monism on an experiential base is not obviously empty:
the foundation can be made non-arbitrary, the geometry and dimension can be forced rather than
chosen, and a wide span of physics and a structural account of mind can be made to come out of
the one rule, in a finite simulation, with results that could have failed and sometimes did
before being corrected. That is a reason to take the picture seriously as a hypothesis. It is
not a reason to take it as established. The distinctive predictions meeting data, and the felt
interior being modeled as well as the shared physical world, are the work that would move it from
a serious hypothesis toward a confirmed theory.
